% ---------------------------------------------------------------------------
% Main theory results, stated in derivation order.
%
% This section is \input{} from theory_appendix.tex right after the roadmap.
% It is the "slow, complete" counterpart of the narrative main text: every
% statement that the paper relies on is written out once, in the order in
% which it is derived, with a pointer to the section that proves it.  The
% narrative main text promotes only a subset and rewrites them for the story.
%
% Labels in this file are prefixed mr- (main results) to avoid collisions.
% ---------------------------------------------------------------------------
\section{Main theory results in derivation order}
\label{sec:app-main-results}

This section collects every theoretical statement used in the paper, in the
order in which it is derived, so that a reader who prefers a linear
exposition can read it as a self-contained summary before the proofs.  Each
statement names the section that proves it.  All results concern the linear
teacher--student model of Section~\ref{sec:mr-setting}; only
Section~\ref{sec:mr-nonlinear} concerns nonlinear feature maps, and it is a
diagnostic rather than a theorem.

\subsection{Setting, estimator and evaluation distributions}
\label{sec:mr-setting}

\paragraph{Data and teacher.}
Inputs are $\x\sim\mathcal N(0,\Sigma)$ with $\Sigma\succeq0$ diagonalized as
$\Sigma=\sum_{k=1}^{d}s_k\uvec_k\uvec_k^\top$.  The teacher is linear with
weight $\bstar=\sum_k b_k\uvec_k$ and the observed response is
\begin{equation}
y=\x^\top\bstar+\sigma\epsilon,\qquad \epsilon\sim\mathcal N(0,1),
\qquad S:={\bstar}^\top\Sigma\bstar=\sum_k s_kb_k^2 .
\label{eq:mr-teacher}
\end{equation}
Given $n$ samples stacked as $X\in\R^{n\times d}$, $\mathbf y=X\bstar+\sigma\eta$,
write $\hat\Sigma=X^\top X/n$ and $\gamma=d/n$.

\paragraph{Students.}
A pixel-space student is $f(\x)=\x^\top\bhat$.  A feature-space student uses a
fixed linear map $F\in\R^{p\times d}$ and reads out $f(\x)=\hat w^\top F\x$;
its effective pixel weight is $\bhat:=F^\top\hat w$, which is also the input
gradient $\nabla_\x f$.  Ridge regression with penalty $\lambda$ gives
\begin{equation}
\bhat=(\hat\Sigma+\lambda I)^{-1}\hat\Sigma\bstar
+\sigma(\hat\Sigma+\lambda I)^{-1}\tfrac1nX^\top\eta
\quad\text{(pixel)},\qquad
\hat w=(\hat\Sigma_F+\lambda I)^{-1}\tfrac1n(FX^\top)\mathbf y
\quad\text{(feature)},
\label{eq:mr-ridge}
\end{equation}
with $\hat\Sigma_F=FX^\top XF^\top/n$.  Software that minimizes
$\|\mathbf y-X\beta\|^2+\alpha\|\beta\|^2$ uses $\alpha=n\lambda$.

\paragraph{Three evaluation distributions.}
The same fitted weight is evaluated on three input distributions.
\begin{enumerate}
\item \emph{Generalization}: $\x\sim\mathcal N(0,\Sigma)$.
\item \emph{Self-accentuation}: $\x=\x_0+\alpha\bhat$, $\alpha\sim p(\alpha)$,
the idealized gradient-ascent path of the student.  The main text takes
$\x_0=0$, $\E[\alpha]=0$, and the variance-matched range
$\Var[\alpha]\,({\bhat}^\top\bhat)^2=S$, so that the student-predicted
response variance along the path equals the natural response variance.
\item \emph{Peer review}: $\x=\x_0+\alpha\bprime$ for a path generated by
another student $\bprime$, with the analogous variance matching.
\end{enumerate}
Nonzero and multiple seeds are treated in
Section~\ref{sec:app-evaluation-distributions}.

\subsection{Metrics on an arbitrary evaluation distribution}

\begin{lem}[Metrics of a linear predictor against a noiseless linear teacher]
\label{lem:mr-metrics}
Let the evaluation inputs have mean $\mu_p$ and covariance $\Sigma_p$, and let
$\Delta:=\bhat-\bstar$.  Then
\begin{align}
\mathrm{MSE}&=\Delta^\top(\mu_p\mu_p^\top+\Sigma_p)\Delta, &
R^2&=1-\frac{\Delta^\top(\mu_p\mu_p^\top+\Sigma_p)\Delta}{{\bstar}^\top\Sigma_p\bstar},\\
\mathrm{slope}&=\frac{{\bhat}^\top\Sigma_p\bstar}{{\bhat}^\top\Sigma_p\bhat}, &
r&=\frac{{\bhat}^\top\Sigma_p\bstar}
{\sqrt{({\bhat}^\top\Sigma_p\bhat)({\bstar}^\top\Sigma_p\bstar)}},
\end{align}
where slope is the regression coefficient of the true response on the
predicted response and $R^2$ is the direct (uncalibrated) coefficient of
determination.  Proof: Lemmas~\ref{lem:pred_error}--\ref{lem:linear_slope}.
\end{lem}

\begin{cor}[The three evaluations]
\label{cor:mr-three-evaluations}
Substituting the three covariances $\Sigma$, $\Var[\alpha]\bhat{\bhat}^\top$,
and $\Var[\alpha]\bprime{\bprime}^\top$ into Lemma~\ref{lem:mr-metrics} gives
Table~\ref{tab:theory-Evaluation-metrics}.  In particular, with
$N:={\bhat}^\top\bstar$, $D:={\bhat}^\top\bhat$ and variance matching,
\begin{equation}
\frac{E_{gen}}{S}=\frac{\Delta^\top\Sigma\Delta}{S},\qquad
\frac{E_{acc}}{S}=\Big(1-\frac ND\Big)^2,\qquad
R^2_{acc}=1-\Big(\frac DN-1\Big)^2,\qquad
\mathrm{slope}_{acc}=\frac ND .
\label{eq:mr-acc-metrics}
\end{equation}
Pearson $r$ is identically $1$ (or undefined) on a noiseless rank-one path and
is therefore not informative for accentuation.  All accentuation metrics are
functions of the single scalar $z_{acc}:=D/N-1$.
Proof: Section~\ref{sec:app-evaluation-distributions}.
\end{cor}

\begin{prop}[Geometry of equal-error sets]
\label{prop:mr-iso-geometry}
Fix $\bstar$ and $\Sigma\succ0$.  In the space of student weights $\bhat$:
\begin{enumerate}
\item the equal-$E_{gen}$ sets are ellipsoids
$\Delta^\top\Sigma\Delta=c$ centered at $\bstar$, with semiaxis
$\sqrt{c/s_k}$ along $\uvec_k$;
\item the equal-$z_{acc}$ sets (hence equal $E_{acc}$, $R^2_{acc}$, slope)
are Euclidean spheres with diameter from $0$ to $(1+z)\bstar$:
$\|\bhat-\tfrac{1+z}{2}\bstar\|^2=\tfrac{(1+z)^2}{4}\|\bstar\|^2$;
\item for a fixed peer $\bprime$, the equal-peer-error sets are hyperplanes
normal to $\bprime$ through $a\bstar$ with $a={\bprime}^\top\bhat/{\bprime}^\top\bstar$.
\end{enumerate}
Proof: Section~\ref{sec:app-iso-error-geometry}.
\end{prop}

\begin{example}[Exact dissociation without learning]
\label{ex:mr-dissociation}
Let $\Sigma=\operatorname{diag}(1,\varepsilon)$, $\bstar=\bhat_1=(1,0)^\top$
and $\bhat_2=(1,\varepsilon^{-q})^\top$ with $0<q<1/2$.  As
$\varepsilon\to0$, $E_{gen}(\bhat_2)=\varepsilon^{1-2q}\to0$ while
$\mathrm{slope}_{acc}(\bhat_2)\to0$ and $R^2_{acc}(\bhat_2)\to-\infty$; both
peer-review scores $R^2_{peer}(\bhat_2\leftarrow\bhat_1)$ and
$R^2_{peer}(\bhat_1\leftarrow\bhat_2)$ equal $1$.  Thus indistinguishable
prediction, arbitrarily poor control, and perfect mutual peer review coexist.
\end{example}

\subsection{Deterministic-equivalent toolkit}
\label{sec:mr-toolkit}

\paragraph{Effective regularization and spectral sums.}
For a population spectrum $\{s_k\}$ and aspect ratio $\gamma=d/n$, let
$\kappa=\kappa(\lambda)$ be the unique positive solution of
\begin{equation}
\kappa-\lambda=\frac{\kappa}{n}\operatorname{Tr}\!\big[\Sigma(\Sigma+\kappa I)^{-1}\big]
=\frac{\kappa}{n}\,\mathrm{df}_1(\kappa),
\label{eq:mr-kappa}
\end{equation}
and define, for $p,q\ge0$,
\begin{equation}
B_{p,q}(\kappa)=\sum_k\frac{s_k^p}{(s_k+\kappa)^q}b_k^2,\qquad
\mathrm{df}_{p,q}(\kappa)=\sum_k\frac{s_k^p}{(s_k+\kappa)^q},\qquad
\mathrm{df}_1:=\mathrm{df}_{1,1},\ \mathrm{df}_2:=\mathrm{df}_{2,2}.
\label{eq:mr-spectral-sums}
\end{equation}
The relation $n\lambda=\kappa\,(n-\mathrm{df}_1)$ and the derivative
$d\kappa/d\lambda=n/(n-\mathrm{df}_2)$ follow from Eq.~\eqref{eq:mr-kappa}.
The notation $A\asymp B$ means that $A-B\to0$ almost surely in the
proportional limit $n,d\to\infty$, $d/n\to\gamma$; the one- and two-resolvent
deterministic equivalents used throughout are listed in
Section~\ref{sec:app-de-toolkit}.

\subsection{Pixel-space ridge regression}

\begin{prop}[Per-eigenmode weight error]
\label{prop:mr-per-pc}
For pixel ridge under Gaussian design,
\begin{equation}
\E\big[(\uvec_k^\top(\bhat-\bstar))^2\big]\asymp
\underbrace{\frac{\kappa^2}{(s_k+\kappa)^2}b_k^2}_{\text{shrinkage bias}}
+\underbrace{\frac{s_k}{(s_k+\kappa)^2}\,
\frac{\kappa^2B_{1,2}(\kappa)+\sigma^2}{n-\mathrm{df}_2(\kappa)}}_{\text{finite-sample signal and noise variance}} .
\label{eq:mr-per-pc}
\end{equation}
The variance term is largest for modes with $s_k\approx\kappa$ and, relative
to the mode's own variance $s_k$, is largest for low-variance modes.
Proof: Section~\ref{sec:app-pixel-ridge}.
\end{prop}

\begin{prop}[Pixel-ridge deterministic equivalents for prediction and control]
\label{prop:mr-pixel-de}
With $\kappa$ from Eq.~\eqref{eq:mr-kappa} on the spectrum of $\Sigma$,
\begin{align}
E_{gen}&\asymp E_{gen,DE}:=\frac{n\big(\kappa^2B_{1,2}+\sigma^2\big)}{n-\mathrm{df}_2}-\sigma^2,
\label{eq:mr-pixel-egen}\\
N&\asymp\mu_N:=B_{1,1},
\label{eq:mr-pixel-N}\\
D&\asymp\mu_D:=B_{2,2}+\frac{\mathrm{df}_{1,2}}{n}\big(E_{gen,DE}+\sigma^2\big).
\label{eq:mr-pixel-D}
\end{align}
Equation~\eqref{eq:mr-pixel-egen} is the sum of Eq.~\eqref{eq:mr-per-pc}
weighted by $s_k$; Eq.~\eqref{eq:mr-pixel-D} is its unweighted sum plus the
squared mean.  The leading (ratio-of-DE) accentuation metrics follow from
Eq.~\eqref{eq:mr-acc-metrics} with $N/D\approx\mu_N/\mu_D$, i.e.
\begin{equation}
z^{(0)}_{acc}=\frac{\mu_D}{\mu_N}-1
=\frac{B_{2,2}-B_{1,1}}{B_{1,1}}
+\frac{\mathrm{df}_{1,2}\,(E_{gen,DE}+\sigma^2)}{n\,B_{1,1}} .
\label{eq:mr-pixel-z}
\end{equation}
The first term is a shrinkage bias ($B_{2,2}\le B_{1,1}$, so it is negative:
ridge shrinkage makes the student under-shoot along its own path); the
second term is Euclidean estimation variance, which prediction weights by
$s_k$ but control counts in full.
Proof: Section~\ref{sec:app-pixel-ridge}.
\end{prop}

\begin{prop}[Prediction-optimal and control-optimal regularization]
\label{prop:mr-pixel-tuning}
\begin{enumerate}
\item If $E_{gen,DE}(\kappa)$ has an interior differentiable minimizer
$\kappa^\star_{gen}$ on the admissible $\lambda\ge0$ path, then
\begin{equation}
E_{gen,DE}+\sigma^2=n\kappa\frac{B_{2,3}}{\mathrm{df}_{2,3}}
\quad\text{at }\kappa=\kappa^\star_{gen}.
\label{eq:mr-pred-optimal}
\end{equation}
\item At a point with $\mu_N>0$, the leading accentuation error vanishes
($\mu_D=\mu_N$, i.e. $z^{(0)}_{acc}=0$, $R^{2,(0)}_{acc}=1$,
$\mathrm{slope}^{(0)}_{acc}=1$) if and only if
\begin{equation}
E_{gen,DE}+\sigma^2=n\kappa\frac{B_{1,2}}{\mathrm{df}_{1,2}}
\qquad\Longleftrightarrow\qquad
\sigma^2=n\lambda\frac{B_{1,2}}{\mathrm{df}_{1,2}} .
\label{eq:mr-control-optimal}
\end{equation}
At such a root $\kappa^\star_{acc}$ the prediction risk is
$E_{gen,DE}(\kappa^\star_{acc})=\sigma^2\,\mathrm{df}_1/(n-\mathrm{df}_1)$.
\item At an interior prediction optimum the residual leading control error is
\begin{equation}
z^{(0)}_{acc,CV}
=\frac{\kappa\,\mathrm{df}_{1,2}}{B_{1,1}}
\left(\frac{B_{2,3}}{\mathrm{df}_{2,3}}-\frac{B_{1,2}}{\mathrm{df}_{1,2}}\right)_{\kappa=\kappa^\star_{gen}},
\label{eq:mr-zacc-cv}
\end{equation}
which is zero only when the two teacher-weighted spectral averages coincide,
e.g. for isotropic $\Sigma$, but not in general for anisotropic spectra.
\end{enumerate}
The zero-error statement concerns the leading ratio-of-DE approximation; the
finite-sample expectation retains $\Var(N/D)$.  Existence, uniqueness, and
boundary optima are discussed in Section~\ref{sec:app-optimal-regularization}.
Proof: Section~\ref{sec:app-optimal-regularization}.
\end{prop}

\begin{prop}[Peer review between independent ridge fits]
\label{prop:mr-peer}
Let $\bhat$ and $\bprime$ be ridge fits from independent samples with the
same $\kappa$.  Then ${\bhat}^\top\bprime\asymp B_{2,2}$ and
${\bstar}^\top\bprime\asymp B_{1,1}$, so the leading peer gain and score are
\begin{equation}
G_{peer}:=\frac{{\bhat}^\top\bprime}{{\bstar}^\top\bprime}\asymp\frac{B_{2,2}}{B_{1,1}},
\qquad
R^{2,(0)}_{peer}=1-\Big(1-\frac{B_{2,2}}{B_{1,1}}\Big)^2 .
\label{eq:mr-peer}
\end{equation}
The peer path removes the Euclidean variance term of
Eq.~\eqref{eq:mr-pixel-D} because independent noise realizations are
uncorrelated; peer review therefore measures shared shrinkage bias only.
Proof: Section~\ref{sec:app-ratio-corrections}, peer-review subsection.
\end{prop}

\subsection{Beyond ratio of means}

\begin{prop}[Second-order ratio correction]
\label{prop:mr-ratio-delta}
Write $N=\mu_N+\delta_N$, $D=\mu_D+\delta_D$.  To second order,
\begin{equation}
\E\Big[\frac ND\Big]\approx\frac{\mu_N}{\mu_D}
\Big(1-\frac{\Cov(N,D)}{\mu_N\mu_D}+\frac{\Var(D)}{\mu_D^2}\Big),\qquad
\E\Big[\Big(1-\frac ND\Big)^2\Big]
=\Big(1-\E\Big[\frac ND\Big]\Big)^2+\Var\Big(\frac ND\Big).
\end{equation}
With $\bhat=g+h$ ($g$ the shrunken signal, $h$ the noise part), the
response-noise contribution is
\begin{equation}
\Var\Big(\frac ND\Big)\approx\frac{1}{\mu_D^2}
\Big[\frac{\sigma^2}{n}\sum_k\frac{s_kb_k^2}{(s_k+\kappa)^2}
\Big(1-\frac{2\mu_N}{\mu_D}\frac{s_k}{s_k+\kappa}\Big)^2
+\frac{2\mu_N^2\sigma^4}{\mu_D^2n^2}\sum_k\frac{s_k^2}{(s_k+\kappa)^4}\Big],
\end{equation}
which omits random-design fluctuations and therefore under-estimates the
variance near ridgeless operating points.  The variance term dominates
$E_{acc}$ precisely when the mean calibration error is small
(e.g. near $\kappa^\star_{acc}$), which is why leading DEs predict the
location of the control optimum but not the value of $E_{acc}$ there.
Proof: Section~\ref{sec:app-ratio-delta}.
\end{prop}

\begin{remark}[Gaussian distributional surrogate]
\label{rem:mr-gaussian-surrogate}
Approximating $\bhat$ by a Gaussian with the DE mean
$\Sigma(\Sigma+\kappa I)^{-1}\bstar$ and diagonal per-mode variance from
Eq.~\eqref{eq:mr-per-pc}, and pushing it through $N$ and $D$, yields
medians, quantiles and the probability of sign changes of $N$.  It is a
surrogate, not a proved distributional deterministic equivalent: it does not
enforce the exact ridgeless identity $N=D$, and it inherits the missing
random-design covariance.  Section~\ref{sec:app-ratio-gaussian}.
\end{remark}

\subsection{Linear feature-space ridge regression}

\begin{lem}[Two population-optimal feature weights]
\label{lem:mr-feature-weights}
Let $C:=\Sigma_F=F\Sigma F^\top$ and $G:=FF^\top$.  The feature weight
minimizing population prediction error and its irreducible error are
\begin{equation}
w^\star=C^\dagger F\Sigma\bstar,\qquad
S_\perp:=\E_\x\big[(\x^\top\bstar-\x^\top F^\top w^\star)^2\big]
=S-{w^\star}^\top Cw^\star ,
\end{equation}
whereas the feature weight whose pixel image best reconstructs the teacher is
$w^P=G^\dagger F\bstar$.  They coincide when $\Sigma\propto I$ on the row
space of $F$ and differ otherwise.  Proof: Section~\ref{sec:app-feature-regression}.
\end{lem}

\begin{prop}[Feature-ridge deterministic equivalents]
\label{prop:mr-feature-de}
Let $\kappa$ solve Eq.~\eqref{eq:mr-kappa} on the spectrum of $C$ with
aspect ratio $p/n$, and define the feature-space spectral sums
\begin{equation}
B^F_{a,b}(\kappa):={w^\star}^\top C^a(C+\kappa I)^{-b}w^\star,\qquad
\mathrm{df}^F_{a,b}(\kappa):=\operatorname{Tr}\!\big[C^a(C+\kappa I)^{-b}\big].
\label{eq:feature-spectral-shorthand}
\end{equation}
Then
\begin{align}
E_{gen,DE}&=\frac{n}{n-\mathrm{df}^F_2}\big(\kappa^2B^F_{1,2}+S_\perp+\sigma^2\big)-\sigma^2,
\label{eq:mr-feature-egen}\\
N_{DE}&={\bstar}^\top\tilde\beta_\kappa,\qquad
\tilde\beta_\kappa:=F^\top(C+\kappa I)^{-1}F\Sigma\bstar
=F^\top C(C+\kappa I)^{-1}w^\star,
\label{eq:mr-feature-N}\\
D_{DE}&=\|\tilde\beta_\kappa\|^2+\frac{E_{gen,DE}+\sigma^2}{n}\,T_F(\kappa),
\qquad
T_F(\kappa):=\operatorname{Tr}\!\big[G(C+\kappa I)^{-2}C\big].
\label{eq:mr-feature-D}
\end{align}
Setting $F=I$ recovers Proposition~\ref{prop:mr-pixel-de}.  The effective
training noise for learning $w^\star$ is $\sigma^2+S_\perp$.
Proof: Section~\ref{sec:app-feature-regression}.
\end{prop}

\begin{cor}[Distinct summary statistics for prediction and control]
\label{cor:mr-summary-statistics}
At leading order, $E_{gen,DE}(\kappa)$ for every $\kappa$, hence the
DE-predicted CV choice $\kappa_{CV}$ and $R^2_{gen}$, are functions of
\begin{equation}
\mathcal P:=\big(C,\ F\Sigma\bstar,\ S,\ \sigma^2,\ n\big)
\end{equation}
only, whereas $N_{DE}$, $D_{DE}$ and every accentuation metric additionally
depend on the Euclidean geometry $G=FF^\top$ (through $T_F$ and
$\|\tilde\beta_\kappa\|^2$) and on the teacher image $F\bstar$ (through
$N_{DE}$).  Fixing $\mathcal P$ fixes prediction; it does not fix control
(Figure~\ref{fig:summary_stats_dep_graph}).
\end{cor}

\begin{prop}[Prediction- and control-optimal feature ridge]
\label{prop:mr-feature-tuning}
An interior minimizer of Eq.~\eqref{eq:mr-feature-egen} satisfies
$E_{gen,DE}+\sigma^2=n\kappa B^F_{2,3}/\mathrm{df}^F_{2,3}$, the same form as
Eq.~\eqref{eq:mr-pred-optimal} with $\sigma^2\mapsto\sigma^2+S_\perp$ inside
$E_{gen,DE}$.  The leading accentuation error vanishes exactly when
\begin{equation}
E_{gen,DE}+\sigma^2
=n\,\frac{{\bstar}^\top\tilde\beta_\kappa-\|\tilde\beta_\kappa\|^2}{T_F(\kappa)}
=n\,\frac{(w^P-w^\star)^\top G\tilde w_\kappa+\kappa\,{w^\star}^\top(C+\kappa I)^{-1}G\tilde w_\kappa}{T_F(\kappa)},
\label{eq:feature-control-optimal}
\end{equation}
with $\tilde w_\kappa:=C(C+\kappa I)^{-1}w^\star$.  The first condition
depends on $\mathcal P$ alone; the second depends on $G$ and $F\bstar$.  An
admissible root is not guaranteed for an arbitrary $F$; in particular, when
$w^P\ne w^\star$ the numerator can be negative for all $\kappa$, in which
case no ridge level calibrates control.
Proof: Section~\ref{app:regularization-derivations}.
\end{prop}

\begin{prop}[Prediction-preserving gauge family]
\label{prop:mr-gauge}
Assume $\Sigma\succ0$ and let
$\mathcal G_b:=\{Q\in O(d):Q\Sigma^{1/2}\bstar=\Sigma^{1/2}\bstar\}$.  For
$Q\in\mathcal G_b$ the map $F\mapsto F_Q:=F\Sigma^{1/2}Q\Sigma^{-1/2}$
preserves $C$ and $F\Sigma\bstar$, hence every element of $\mathcal P$, the
Gaussian training law of $(F\x,y)$, the full curve $E_{gen,DE}(\kappa)$ and
the DE-based CV choice; it changes $G$ and $F\bstar$ in general, and
therefore changes $N_{DE}$, $D_{DE}$ and all accentuation metrics.  The
subgroup $\mathcal G_{b,c}:=\{Q\in\mathcal G_b:Q\Sigma^{-1/2}\bstar=\Sigma^{-1/2}\bstar\}$
additionally fixes $F\bstar$, isolating $G$ as the only control-relevant
statistic that varies.  When $\Sigma=I$ the family acts trivially on both
prediction and control.  Members of the family need not give identical CV
choices for the same finite realization of $X$.
Proof: Section~\ref{sec:app-gauge-fragility}.
\end{prop}

\begin{cor}[Continuous prediction-null rotation]
\label{cor:mr-rotation}
Write $A:=F\Sigma^{1/2}$ and let $P$ be the part of $A$ needed to fix
$A\Sigma^{1/2}\bstar$ (and optionally $A\Sigma^{-1/2}\bstar$).  For any two
row frames $R,V$ in the orthogonal complement with $RR^\top=VV^\top$ and
$RV^\top=0$, the path $A(t)=P+\sqrt{1-t}\,R+\sqrt t\,V$, $F(t)=A(t)\Sigma^{-1/2}$,
$t\in[0,1]$, lies in the gauge family.  Choosing $R$ in high-variance and $V$
in low-variance pixel directions moves feature energy continuously from the
spectral head to the tail: $E_{gen}$, $R^2_{gen}$ and $\kappa_{CV}$ are
constant in $t$ while $\operatorname{Tr}(F(t)F(t)^\top)$, $T_F$ and
$z_{acc}$ grow without bound.
Proof: Section~\ref{sec:app-gauge-fragility}, construction subsection.
\end{cor}

\begin{prop}[Spectral form of the fragility factor]
\label{prop:mr-fragility}
Let $C=\sum_ks_kv_kv_k^\top$ and $g_k:=F^\top v_k$ be the pixel gradient of
feature PC $k$.  Then
\begin{equation}
T_F(\kappa)=\sum_{k:s_k>0}\Big(\frac{s_k}{s_k+\kappa}\Big)^2
\frac{\|g_k\|^2}{g_k^\top\Sigma g_k},
\qquad s_k=g_k^\top\Sigma g_k .
\label{eq:mr-fragility}
\end{equation}
$T_F$ is a ridge-weighted average, over feature PCs that survive shrinkage,
of the ratio between the Euclidean and the natural-covariance energy of the
backpropagated gradient.  It is response-independent and large when feature
PCs backpropagate into low-variance (high spatial frequency) input
directions.  Proof: Section~\ref{sec:app-feature-fragility}.
\end{prop}

\begin{remark}[When is control fragile although prediction is fine]
\label{rem:mr-when}
Combining Propositions~\ref{prop:mr-feature-de}, \ref{prop:mr-feature-tuning}
and \ref{prop:mr-fragility}, the leading control error at fixed prediction
error grows with three separable factors:
(i) the estimation-noise level $(E_{gen,DE}+\sigma^2)/n$;
(ii) the spectral mismatch of Eq.~\eqref{eq:mr-zacc-cv}, i.e. how far the
CV-selected $\kappa$ sits from the control-calibrated $\kappa$, which
vanishes for isotropic data and grows with anisotropy and with teacher
energy in low-variance modes;
(iii) the gradient-geometry factor $T_F$, which vanishes only if surviving
feature gradients live in high-variance input directions.
Only factor (i) is visible from held-out prediction.
\end{remark}

\subsection{Nonlinear feature maps: a diagnostic, not a theorem}
\label{sec:mr-nonlinear}

\paragraph{Local diagnostic.}
For a differentiable feature map $\phi:\R^d\to\R^p$ with Jacobian
$J_\phi(x)$ and feature covariance $\Sigma_\phi=\sum_ks_kv_kv_k^\top$ on
natural inputs, define
\begin{equation}
T_\phi(x,\kappa):=\sum_k\frac{s_k}{(s_k+\kappa)^2}\|J_\phi(x)^\top v_k\|^2,
\qquad
V_\phi:=\frac{E_{val}}{n}\,\E_x[T_\phi(x,\kappa)],
\label{eq:mr-nonlinear-T}
\end{equation}
where $E_{val}$ is the held-out prediction error and $\kappa$ is computed
from the empirical feature spectrum and the CV-selected penalty.  For
$\phi(x)=Fx$ this equals $T_F$ and the variance term of
Eq.~\eqref{eq:mr-feature-D}.  For nonlinear $\phi$ the identity
$s_k=g_k^\top\Sigma g_k$ fails, so Eq.~\eqref{eq:mr-nonlinear-T} is a
diagnostic motivated by the linear theory; it is evaluated in exact-local,
finite-difference, smoothed-Jacobian, neighborhood and finite-step versions
(Section~\ref{sec:app-nonlinear-diagnostic}).

\paragraph{Empirical claim.}
On the neuronal control data (25 sites, 10 encoding models), after removing
site means, $V_\phi$ and $T_\phi$ are associated with control slope and with
direct control error across the nine trained models, whereas held-out
prediction error is not; the association is weaker but same-signed among the
seven conventionally trained models.  These are descriptive correlations with
repeated models within sites, not independent-sample tests
(Section~\ref{sec:app-nonlinear-biological-validation}).
